% ===================================================================== RELATED WORK
\section{Related work}
\label{sec:related}

\paragraph{Quorums, consensus, and reconfiguration}
Deterministic asynchronous consensus faces the FLP limit~\cite{flp}; practical BFT protocols
typically rely on partial synchrony~\cite{dls,pbft}. Byzantine quorum systems formalize the
intersection and availability conditions used here~\cite{byzantinequorums}. Flexible and
multi-threshold BFT already permit different decision types to use different quorums
\cite{flexiblebft,multithresholdbft}. Thus neither our boundary threshold nor our fence-only
inequality is new quorum mathematics. The distinction is where those conditions apply in an
engine handoff, and what durable actions their signers have taken.

Potential engines include HotStuff~\cite{hotstuff}, HotStuff-1~\cite{hotstuff1},
Tendermint~\cite{tendermint}, DAG-based protocols
\cite{narwhal,bullshark,shoal,shoalpp,autobahn,mysticeti}, and the broadcast/consensus hybrid
Sui Lutris~\cite{suilutris}. Sui Lutris proves reconfiguration of its persistent hybrid; it is
not a protocol for retiring one authoritative engine and bootstrapping a different engine across
fault models. Membership reconfiguration work addresses changing replica sets
\cite{verticalpaxos,reconfig,bftsmart,mongoraft,dynanet,hydra}. Stoppable Paxos ends a
configuration by a stop command before its successor begins~\cite{stoppablepaxos}; asynchronous
Byzantine reconfiguration addresses decisions by a superseded configuration
\cite{asyncreconfig}. These works motivate our source-generation retirement obligation, but
\sage{} holds committee membership fixed and changes the engine instead. A HotStuff-1 target,
for example, would need an adapter that distinguishes its speculative confirmation from the
irreversible finality represented in our two-tier contract; no such adapter is evaluated here.

\paragraph{Deployed transitions}
The Ethereum Merge prepared a proof-of-stake chain in parallel and used terminal total difficulty
to trigger the transition~\cite{ethereummerge,eip3675}. EIP-3675 allows more than one terminal
proof-of-work block, so the trigger alone need not name a unique checkpoint. Gasper, GRANDPA,
Tendermint, and Tenderbake are consensus or finality mechanisms rather than specifications for
this fixed-committee engine handoff~\cite{gasper,grandpa,tendermint,tenderbake}.
Hyperledger Fabric~v3 supplies a closer \emph{operational} comparison: its Raft-to-BFT procedure
keeps the consenter set unchanged and configures the BFT service for $n=3f+1$
\cite{fabricbftmigration}. The procedure enters maintenance mode, rejects normal transactions,
stops and restarts ordering nodes, and cannot return to Raft after BFT begins committing. Its
ordered configuration and maintenance interval delimit the transition without a
\sage{}-style provisional suffix. The design difference is therefore not ``halt versus no halt'':
\sage{} also pauses finalization for certification, while keeping source service during preparation
and specifying a certificate-gated abort before seal. Neither procedure was measured against the
other in a shared deployment.

\paragraph{Runtime and adaptive switching}
Run-time switching between total-order algorithms predates blockchains~\cite{runtimeorder}.
Abstract composes BFT instances using idempotent requests and signed initialization
histories~\cite{abstracteurosys,abstract}. Subsequent systems include a Zab-to-Instant-Commit
ZooKeeper switch~\cite{zkswitch} and Byzantine switching algorithms in
BFT-SMaRt~\cite{dynamicbft}. BFTBrain selects among six BFT protocols in one cluster and uses an
Abstract-derived handoff wrapper~\cite{bftbrain}. Cox uses a committed configuration block,
matching signed termination checkpoints, and lagging-node recovery to switch among BFT cores;
its security argument permits different core quorum values~\cite{cox}. ATBFT switches between
leader-based and leaderless BFT modes~\cite{atbft}. Other frameworks adapt protocol choice to
conditions~\cite{spectrum,adapt,aware,dapbft,adachain,adaptbft,threatadapt,adacon}, whereas
resilient designs harden a protocol family~\cite{rebft,prime,aardvark,rbft}. Bedrock provides
a common BFT implementation substrate, not itself an engine-migration protocol~\cite{bedrock}.

Quorum-family change alone cannot establish novelty: Cox addresses unequal BFT-core quorums,
Abstract permits distinct instance clusters, and Fabric crosses from crash to Byzantine fault
tolerance. For the fixed, identical committee studied here, \sage{} instead specifies the
dependency between a durable readiness lock and source quiescence, policy-bound retirement
evidence, complete-manifest identity, and a certified provisional suffix. The mechanisms in
\cref{tab:proof-obligations,tab:switchcompare} are source-described, not a head-to-head performance
result; none of the cited implementations was run in this study. BFTBrain's same-cluster
optimization starts the next epoch from already executed state without waiting again for
matching initialization histories; its Appendix~B explains how the Abstract wrapper preserves
composition~\cite{bftbrain,abstracteurosys}.

\begin{table*}[t]
\centering
\caption{Source-described handoff obligations. ``Not specified'' identifies the reviewed source's
boundary, not a demonstrated vulnerability. No row is a measured comparison.}
\label{tab:proof-obligations}
\small
\setlength{\tabcolsep}{3pt}
\renewcommand{\arraystretch}{1.14}
\begin{tabularx}{\textwidth}{@{}L{0.13\textwidth}YYYY@{}}
\toprule
\textbf{Obligation} & \textbf{Abstract / BFTBrain} & \textbf{Cox} & \textbf{Fabric v3} & \textbf{\sage{} design} \\
\midrule
Boundary uniqueness & Valid signed init history; BFTBrain commits an epoch-ending request & Committed config block and matching signed checkpoints & Ordered channel configuration; maintenance boundary & One durable readiness body per epoch; $n-f$ \CutCert \\
Old-authority exclusion & Backup aborts later requests; composition relies on each instance's correctness & Reconfiguring core rejects blocks above config block; checkpoint precedes replacement & Normal transactions rejected; ordering nodes stopped and restarted & Durable Rule~\ref{rule:lock} and $q_1>2f$; separate policy-bound retirement evidence \\
State bootstrap & Execute init history; same-cluster BFTBrain already holds the state & Termination status binds next-core config and consensus state; genesis initializes new core & In-place ledger and BFT configuration; quorum checked before service resumes & Certified parent and root; adapter supplies native metadata and safety predicate \\
Rollback semantics & Request abort and retry, not reversion of committed history & Lagging-epoch recovery; no reversible committed suffix specified & Backup recovery before exit from maintenance; one-way after BFT commits & Provisional suffix; certified abort to fresh generation, or certified seal \\
Committee & Abstract allows separate clusters; BFTBrain uses the same cluster & Compatible BFT cores; distinct core quorums permitted & Same consenters; $n=3f+1$ & One fixed identical committee; source policy explicitly admitted \\
Durability & Instance correctness; crash-store protocol not supplied by the wrapper argument & Core/checkpoint state; atomic write/release/restart ordering not specified in protocol pseudocode & Filesystem backup, retained ledgers, restart procedure & Persist-before-release and reload-before-use required; full runtime conformance unestablished \\
\bottomrule
\end{tabularx}
\end{table*}

\paragraph{Comparator reproducibility and shared dimensions}
Cox publishes both the normal switch and lagging-node recovery: a committed configuration block
puts the old core into a reconfiguring state; matching signed checkpoints certify its termination
status; the new core initializes from that state; and chained epoch-change proofs bring lagging
nodes forward~\cite[Sections~4.2--4.3 and 5.2]{cox}. Its argument uses the larger of the old
and new core quorum values. A missing published checkpoint rule is therefore \emph{not} the reason
for using an internal control. Rather, this study has not established a version-pinned executable
Cox reproduction, compatible codec and state-transfer path, or exact durable
write/release/restart ordering. The RQ1 ``Cox-style'' arm implements none of Cox's checkpoint or
recovery path: it is an internal switch-class cost control, not Cox, and cannot support an
empirical comparison with Cox.

BFTBrain's documented predefined-order mode can isolate wrapper switching from learning in a
shared BFT-to-BFT domain, but source inspection is not a successful execution or a matched
benchmark. A common-harness study would fix hardware, network, request and batch sizes, load,
committee, fault budget, and client-completion semantics before comparing throughput, latency,
and boundary interruption. Fabric's maintenance interval is a separate shared dimension;
\sage's provisional-suffix/abort semantics require a distinct report, not a fabricated
competitor score. Neither a common-harness execution nor a technical impossibility of one is
established here.

\begin{table*}[t]
\centering
\caption{Mechanism comparison of runtime switching and engine transition, ordered by year of the
cited source. No prior implementation was executed in this study.}
\label{tab:switchcompare}
\small
\setlength{\tabcolsep}{4pt}
\renewcommand{\arraystretch}{1.12}
\rowcolors{2}{rowtint}{white}
\begin{tabularx}{\textwidth}{@{}L{0.14\textwidth}c L{0.20\textwidth} Y L{0.21\textwidth}@{}}
\toprule
\textbf{System} & \textbf{Year} & \textbf{Setting / switch} & \textbf{Handoff and recovery} & \textbf{Evidence in source} \\
\midrule
Total-order switch~\cite{runtimeorder} & 2006 & Total-order broadcast algorithms & Run-time adaptation; Byzantine faults out of scope & Feasibility demonstration \\
Abstract / Backup~\cite{abstracteurosys,abstract} & 2010 & BFT instance composition & Idempotent requests; signed initialization histories & Composition theorem; evaluation \\
ZooKeeper switch~\cite{zkswitch} & 2017 & Zab $\to$ Instant Commit, crash faults & Adaptation logic; transfer protocol not detailed & Proof-of-concept \\
BFT-SMaRt switch~\cite{dynamicbft} & 2018 & Among Byzantine protocols & Several switching algorithms in one framework & Comparative implementation \\
SeeMoRe~\cite{seemore} & 2020 & One hybrid crash/Byzantine protocol & Mode change inside one protocol; no engine replacement & Design and evaluation \\
Bedrock~\cite{bedrock} & 2024 & Twelve BFT protocols, one platform & Not a switching protocol; common substrate & Uniform re-implementation \\
BFTBrain~\cite{bftbrain} & 2025 & Six protocols, one BFT cluster & Signed init history; post-boundary requests aborted & Safety/liveness argument; public artifact \\
ATBFT~\cite{atbft} & 2025 & Leader-based $\leftrightarrow$ leaderless BFT & Network-driven selection; uninterrupted processing claimed & Implementation and evaluation \\
Cox~\cite{cox} & 2026 & Online switching among BFT algorithms & Recovery sub-protocol for lagging nodes & Partial-synchrony safety/liveness claim \\
Fabric~v3~\cite{fabricbftmigration} & 2026 & Raft (crash fault tolerant) $\to$ SmartBFT, same consenters, $n=3f+1$ & Maintenance-mode halt; irreversible once BFT commits & Shipped operational procedure \\
\midrule
\rowcolor{sageblue!8}
\textbf{\sage{} (this work)} & 2026 & Fixed committee; PoA $\to$ HotStuff across fault models & Certified boundary, manifest, fence, abort, seal; provisional suffix; deadline-bounded rollback & Exclusion theorems; bounded models with falsifying controls; process controls \\
\bottomrule
\end{tabularx}
\end{table*}
